\documentclass[10pt]{article}
\usepackage{fontspec}
\setmainfont{Inter}[
  Path = fonts/,
  Extension = .ttf,
  UprightFont = *-Regular,
  BoldFont = *-SemiBold,
  ItalicFont = *-Italic,
]
\newfontfamily\playfair{PlayfairDisplay}[
  Path = fonts/,
  Extension = .ttf,
  UprightFont = *-Regular,
  BoldFont = *-Bold,
  ItalicFont = *-Italic,
]
\usepackage[letterpaper,margin=0.7in,top=0.55in,bottom=0.5in]{geometry}
\usepackage{xcolor}
\usepackage[colorlinks=false,pdfborder={0 0 0}]{hyperref}
\usepackage[normalem]{ulem}
\usepackage{enumitem}
\usepackage{setspace}

% Same Turquoise Oasis tokens as the CV/website, on a plain white page.
\definecolor{cvink}{RGB}{8,40,31}
\definecolor{cvmuted}{RGB}{62,109,96}
\definecolor{cvaccent}{RGB}{11,111,88}
\definecolor{cvhairline}{RGB}{190,215,208}
\color{cvink}

\setlength{\parindent}{0pt}
\setlength{\parskip}{0pt}
\pagestyle{empty}
\raggedbottom
\setstretch{1.04}

\newcommand{\link}[2]{\href{#1}{\color{cvaccent}\uline{#2}}}

\newcommand{\rsec}[1]{%
  \par\vspace{8pt}{\bfseries\fontsize{8.5}{10}\selectfont\color{cvaccent}\MakeUppercase{#1}}\par\vspace{3pt}%
  \nointerlineskip{\color{cvhairline}\rule{\linewidth}{0.8pt}}\par\vspace{3pt}}

% Title (bold) with right-aligned dates, then italic muted subtitle.
\newcommand{\rjob}[3]{%
  \par\vspace{3pt}\noindent{\bfseries #1}\hfill{\fontsize{9}{11}\selectfont\color{cvmuted}#2}\par
  {\itshape\fontsize{9}{11}\selectfont\color{cvmuted}#3}\par\vspace{1pt}}

\setlist[itemize]{leftmargin=11pt,itemsep=1pt,topsep=1pt,parsep=0pt,partopsep=0pt,
  label={\color{cvaccent}\textbullet}}

% Compact two-column row for skills / languages.
\newcommand{\rrow}[2]{\par\noindent\parbox[t]{1.7in}{\strut\bfseries #1\strut}\parbox[t]{\dimexpr\linewidth-1.7in\relax}{\strut\color{cvmuted}#2\strut}\par\vspace{1.5pt}}

% Bibliography-style entry with hanging indent.
\newcommand{\rpub}[1]{\par\noindent\hangindent=11pt\hangafter=1\raggedright #1\par\vspace{2.5pt}}

\catcode`\_=12\relax

\begin{document}
\fontsize{9.3}{12}\selectfont

{\color{cvaccent}\fontsize{8}{10}\selectfont\bfseries \MakeUppercase{Educational Measurement \textbullet\ Validity \textbullet\ Psychometrics \textbullet\ AI in Assessment}}\par\vspace{4pt}
{\playfair\fontsize{28}{32}\selectfont Sergio Araneda, Ph.D.}\par\vspace{4pt}
{\fontsize{9}{11}\selectfont\color{cvmuted} sondaxius@gmail.com \ \textbar\ \link{https://sndxs.github.io/personalwebsite/}{sndxs.github.io/personalwebsite} \ \textbar\ \link{https://github.com/sndxs}{github.com/sndxs} \ \textbar\ \link{https://scholar.google.com/citations?user=_Q0VXwsAAAAJ}{Google Scholar}}\par\vspace{6pt}

Research scientist in educational measurement and psychometrics (Ph.D., UMass Amherst) working on validity theory, test security, and the use of AI and large language models in assessment. Industry experience at Caveon and College Board; peer-reviewed and preprint research, a Routledge book chapter, and 45+ conference sessions at NCME, ITC, AERA, and COTS.

\rsec{Experience}

\rjob{Research Scientist, Caveon}{July 2023 -- Present (part-time since July 2026)}{Remote, from New York City, NY}
\begin{itemize}
\item Research on randomly parallel testing and on using LLMs with Caveon's SmartItem\texttrademark\ technology for item generation and test security.
\item Author and co-author of preprints on LLM-based item generation, reliability of randomly parallel tests, and AI-related test security risks; regular presenter at NCME, ITC, COTS, and ATP.
\end{itemize}

\rjob{Associate Psychometrician, College Board}{October 2022 -- April 2023}{Remote, from New York City, NY}
\begin{itemize}
\item Member of the PsyTech Team (Psychometrics): research on item parameter drift and on automated essay scoring for the SAT Essay.
\end{itemize}

\rjob{Research Assistant, DEMRE, Universidad de Chile}{August 2017 -- December 2018}{Santiago, Chile, and remote from Jersey City, NJ}
\begin{itemize}
\item Evaluated the use of the ``ranking de notas'' score in the Chilean university admission system (PSU).
\end{itemize}

\rjob{Research Interns: Caveon (2022) and ETS AI Labs (2021)}{Summers}{Remote, from Amherst, MA}
\begin{itemize}
\item Caveon: built SmartItems and used NLP to generate correct/incorrect statements automatically from text.
\item ETS AI Labs: designed feedback logic for filler words and pauses in the ``Authentic'' interview-practice app.
\end{itemize}

\rjob{Earlier career in finance}{2013 -- 2019}{CPQi (Murex consultant for RBC), Murex (Back-Office consultant), Vision Advisors (quant analyst)}

\rsec{Education}
\rjob{Ph.D., Research, Educational Measurement, and Psychometrics}{2019 -- 2022}{University of Massachusetts Amherst}
Dissertation: \textit{An Experiential Approach to Test Design and Validation}.
\rjob{Mathematical Civil Engineer}{2005 -- 2013}{Universidad de Chile}

\rsec{Technical Skills \& Languages}
\rrow{Programming}{Python (since 2013; PyTorch), R (since 2018), SQL, MATLAB, Java, UNIX, Stata, Git}
\rrow{Psychometric software}{flexMIRT, Mplus; optimization: Gurobi, CPLEX}
\rrow{Languages}{Spanish (native); English (advanced; TOEFL iBT 102/120); French (advanced reading)}


\newpage
\rsec{Selected Publications}
\rpub{Sireci, S. G., Araneda, S., \& Su\'arez-\'Alvarez, J. (2026). The impact of artificial intelligence on validity and test validation. In \textit{The Role of AI in Assessment} (pp.\ 312--336). Routledge. \link{https://doi.org/10.4324/9781003621522-13}{doi:10.4324/9781003621522-13}}
\rpub{Araneda, S., Fremer, J., \& Foster, D. (2024). Commentary: Past, present, and future of educational measurement. \textit{Educational Measurement: Issues and Practice}. \link{https://doi.org/10.1111/emip.12632}{doi:10.1111/emip.12632}}
\rpub{Araneda, S., Lee, D., Lewis, J., Sireci, S. G., Moon, J. A., Lehman, B., Arslan, B., \& Keehner, M. (2022). Exploring relationships among test takers' behaviors and performance using response process data. \textit{Education Sciences}. \link{https://doi.org/10.3390/educsci12020104}{doi:10.3390/educsci12020104}}
\rpub{Araneda, S. (2026). Tests as socio-technical artifacts: A taxonomy for invalidity hypotheses [Preprint]. \link{https://doi.org/10.13140/RG.2.2.33888.65285}{doi:10.13140/RG.2.2.33888.65285}}
\rpub{Araneda, S. (2026). The multiple crises of validity theory in the age of AI [Preprint]. \link{https://doi.org/10.13140/RG.2.2.32698.15048}{doi:10.13140/RG.2.2.32698.15048}}
\rpub{Araneda, S., \& Sireci, S. G. (2026). Making the case for an experiential approach to test validation [Preprint]. \link{https://doi.org/10.13140/RG.2.2.28455.44968}{doi:10.13140/RG.2.2.28455.44968}}
\rpub{Foster, C., \& Araneda, S. (2026). Assessing pre-exposure risks of LLM-generated item pools [Preprint]. \link{https://doi.org/10.13140/RG.2.2.25885.37602}{doi:10.13140/RG.2.2.25885.37602}}
\rpub{Araneda, S., Marder, A., Foster, C., \& Foster, D. F. (2025). Reliability in randomly parallel tests: A practical guide [Preprint]. \link{https://doi.org/10.13140/RG.2.2.22375.28329}{doi:10.13140/RG.2.2.22375.28329}}
\rpub{Araneda, S., Beiting-Parrish, M., Bottoms, B., La Torre, D., et al. (2023). Statement about the new edition of the Standards for Educational \& Psychological Testing [White paper]. NCME GSIC. \link{https://doi.org/10.31235/osf.io/dv4pa}{doi:10.31235/osf.io/dv4pa}}
{\fontsize{8.5}{10}\selectfont\color{cvmuted}Full list, including all preprints and newspaper op-eds in \textit{La Tercera}, \textit{El Mercurio}, and \textit{El Desconcierto}: \link{https://sndxs.github.io/personalwebsite/}{sndxs.github.io/personalwebsite}.}\par

\rsec{Selected Presentations \& Sessions}
\begin{itemize}
\item \textbf{ITC 2026, Auckland:} \textit{The Multiple Crises of Validity Theory in the Age of AI}; chaired two symposia (AI in testing and certification; Hispanic-American perspectives on validity).
\item \textbf{NCME 2026, Los Angeles:} \textit{Factor analysis in randomly parallel tests: Limitations and emerging solutions}.
\item \textbf{AIME-Con 2026, Pittsburgh (October 6):} \textit{Tests as Socio-Technical Artifacts: A Taxonomy for Invalidity Hypotheses}; chair of the coordinated session on validating AI in educational testing.
\item \textbf{AIME-Con 2025, Pittsburgh:} training session on generative AI for item construction; organized session on validating AI in educational testing; talks on validity of AI tools and on tests as socio-technical artifacts.
\item \textbf{COTS 2025 and 2024:} organized a session on AI-driven threats to test security; panelist on generative AI for state assessment programs; SDaR risk-metric talk.
\item \textbf{NCME 2023--2025; ATP Innovations in Testing 2024--2025; IMPS 2024; ITC 2024; NERA; AERA:} talks and demonstrations on AI item generation, randomly parallel tests, and test security; organizer of NCME 2025 discussion on 100 years of criticism of standardized testing.
\end{itemize}

\rsec{Service \& Teaching}
\begin{itemize}
\item NCME: Website Committee (2025--26); Diversity Issues in Testing Committee (2024--25); Chair, Graduate Student Issues Committee (2021--23), coordinating the Standards Study Group (28 recorded sessions).
\item FEVED (Chilean Forum of Educational Assessment Professionals): founding coordinator (2020--21), vice-coordinator (2025--26); organizer of two interview series and a 2025 seminar.
\item Teaching: created and taught \textit{Principles of Game Design Applied in Education} (UMass Amherst, Fall 2021); courses in probability and statistics, linear algebra, and mathematical reasoning (Universidad de Chile).
\end{itemize}

\end{document}
